%% Question bank chapter file - generated by bank_tools.py build (do not edit while the bank is still being built:
%% it is regenerated from the per-test data in Question Bank\data). Every question is one bankq block:
%%   \begin{bankq}[id=..., source={...}, paper=.., qno=.., type=..., subtopic={...}, answer={...}, answer_src=key|solved,
%%                 working={...}, difficulty=1-10, diff_parts={K S I C F}, diff_note={...},
%%                 group=<passage id or empty>, also_in={duplicate copies in other tests}]
%% difficulty = min(10, K+S+I+C+F-1): K concepts 1-3, S steps 1-4, I insight 0-2, C computation 0-2, F format 0-1
%%   ...question LaTeX, options, figures (\bfig / \bfigi -> img/)...
%%   \end{bankq}
%% Passages shared by several questions are bankpassage blocks; their questions carry group=<passage id>.
\documentclass[10pt]{article}
\usepackage[a4paper,margin=1.8cm]{geometry}
\usepackage[T1]{fontenc}
\usepackage{amsmath,amssymb,mathtools}
\usepackage{graphicx}
\usepackage[export]{adjustbox}
\usepackage{xcolor,needspace,keyval}
\usepackage{enumitem}
\usepackage{booktabs,array,tabularx}
\usepackage[hidelinks]{hyperref}
\setlength{\parindent}{0pt}
\setlength{\parskip}{3pt}
\makeatletter
\define@key{bq}{id}{\def\bq@id{#1}}\define@key{bq}{source}{\def\bq@src{#1}}\define@key{bq}{paper}{\def\bq@paper{#1}}
\define@key{bq}{qno}{\def\bq@qno{#1}}\define@key{bq}{type}{\def\bq@type{#1}}\define@key{bq}{subtopic}{\def\bq@sub{#1}}
\define@key{bq}{answer}{\def\bq@ans{#1}}\define@key{bq}{answer_src}{\def\bq@asrc{#1}}\define@key{bq}{working}{\def\bq@work{#1}}
\define@key{bq}{group}{\def\bq@group{#1}}\define@key{bq}{also_in}{\def\bq@also{#1}}
\define@key{bq}{difficulty}{\def\bq@diff{#1}}\define@key{bq}{diff_parts}{\def\bq@dparts{#1}}\define@key{bq}{diff_note}{\def\bq@dnote{#1}}
\newenvironment{bankq}[1][]{\def\bq@work{}\def\bq@also{}\def\bq@group{}\def\bq@diff{?}\def\bq@dparts{}\def\bq@dnote{}\setkeys{bq}{#1}%
  \par\needspace{4\baselineskip}\medskip\noindent{\small\textbf{\bq@id}\hfill\color{gray}[\bq@src, Paper \bq@paper, Q\bq@qno]
  \bq@type\ $\cdot$ \bq@sub\ $\cdot$ \textbf{D\bq@diff}}\par\nopagebreak\relax}
  {\par{\small\color{teal}Answer: \textbf{\bq@ans}\ifx\bq@work\@empty\else\ (solved: \bq@work)\fi
  \ifx\bq@also\@empty\else\par\color{gray}Also in: \bq@also\fi}\par}
\define@key{bp}{id}{\def\bp@id{#1}}\define@key{bp}{source}{\def\bp@src{#1}}\define@key{bp}{questions}{\def\bp@qs{#1}}
\newenvironment{bankpassage}[1][]{\setkeys{bp}{#1}\par\needspace{4\baselineskip}\bigskip\noindent
  {\small\textbf{Passage \bp@id}\hfill\color{gray}[\bp@src; questions \bp@qs]}\par\nopagebreak\relax\itshape}{\par}
\makeatother
\newcommand{\bfig}[2]{\par\smallskip{\centering\includegraphics[width=#1pt,max width=0.62\linewidth,max height=0.32\textheight]{img/#2}\par}\smallskip}
\newcommand{\bfigi}[2]{\includegraphics[width=#1pt,max width=\linewidth,max height=0.32\textheight,valign=c]{img/#2}}
\newcommand{\bmissing}[1]{\textit{[figure missing in the source conversion: \detokenize{#1}]}}
\newenvironment{qoptions}{\begin{itemize}[leftmargin=1.8em,itemsep=1pt,topsep=2pt,label={}]}{\end{itemize}}
\newcommand{\opt}[2]{\item[(#1)] #2}
\providecommand{\fwstart}{}\providecommand{\fwend}{}\providecommand{\qfullsplit}{}
\newcolumntype{L}[1]{>{\raggedright\arraybackslash}p{#1}}
\newcolumntype{C}[1]{>{\centering\arraybackslash}p{#1}}
\newcolumntype{R}[1]{>{\raggedleft\arraybackslash}p{#1}}
\providecommand{\sectionhdr}[1]{}\providecommand{\subsectionhdr}[1]{}


\begin{document}
{\Large\bfseries RBD --- Question Bank}\par\smallskip
{\small Physics chapter P05. Questions from all converted test papers, grouped by source series. Answers: printed key where it matches the paper, otherwise solved (answer\_src=solved, with working). Exact repeats across papers appear once, with also\_in listing the other copies.}\par\medskip\hrule

\section*{Score-I papers}
\begin{bankq}[id=SI12-P1-PHY-S1-Q06, source={Score-I - Test 12 (09 Mar 2022)}, paper=1, qno=6, type=multi, subtopic={Torque}, answer={A,B}, answer_src=key, difficulty=4, diff_parts={K2 S2 I0 C0 F1}, diff_note={concepts: rolling of a spool, torque and force equations | steps: constraint a = alpha r about the contact thread; force equation along the incline; torque equation; solve a and T}]
A spool of mass $m$ and moment of inertia $I$ about its axis rests on a smooth inclined plane, inclined at an angle $\theta$. The end of the thread is tied as shown. The outer radius of the spool is $R$ \& inner radius of the spool is $r$:
\bfig{98.0}{SI12_p1_phy_s1_q06_1.png}
\begin{qoptions}
\opt{A}{Acceleration of centre of mass of spool $a_{cm}=\dfrac{mg\sin\theta}{m+\frac{I}{r^2}}$}
\opt{B}{Tension in thread is $T=\dfrac{mgI\sin\theta}{I+mr^2}$}
\opt{C}{Angular acceleration of spool is $\alpha=\dfrac{a_{cm}}{R}$}
\opt{D}{Acceleration of centre of mass is zero.}
\end{qoptions}
\end{bankq}

\section*{TOAS papers}
\begin{bankq}[id=TOAS03-P2-PHY-S2-Q02, source={TOAS - Test 03 (27 Feb 2022)}, paper=2, qno=2, type=numerical, subtopic={Rotational kinematics}, answer={15.00}, answer_src=key, difficulty=2, diff_parts={K1 S1 I0 C0 F1}, diff_note={concepts: rigid body velocities | steps: angular velocity from the end A; velocity of B}]
The two ends $A$ and $B$ of a uniform rod of length $\ell=1$ m are moving with velocities $V_A$ and $V_B$ as shown. The velocity $V_B$ (m/s) is
\bfig{166.0}{TOAS03_p2_phy_s2_q02_1.png}
\end{bankq}
\begin{bankpassage}[id=TOAS04-P2-PHY-S3-G3-4, source={TOAS - Test 04 (15 Mar 2022)}, questions={3-4}]
A uniformly charged non conducting rod is suspended vertically at its end. The rod can swing freely in the vertical plane without any friction. The linear charge density on the rod is $\lambda$ C/m and it has a uniformly distributed mass of $M$. Length of the rod is $L$. A uniform horizontal electric field ($E$) is switched on in the region.
\bfig{36.2}{TOAS04_p2_phy_s3_para3_4_1.png}
\end{bankpassage}
\begin{bankq}[id=TOAS04-P2-PHY-S3-Q03, source={TOAS - Test 04 (15 Mar 2022)}, paper=2, qno=3, type=numerical, subtopic={Torque}, answer={1.00}, answer_src=key, difficulty=4, diff_parts={K2 S1 I1 C0 F1}, diff_note={concepts: torque, electric force on a charged rod | steps: electric torque about the hinge; energy or torque balance at the horizontal position | I: energy method}, group=TOAS04-P2-PHY-S3-G3-4]
If value of electric field is $E_0$ so that the rod just manages to make itself horizontal. If $E_0=\dfrac{nMg}{\lambda L}$, find $n$.
\end{bankq}
\begin{bankq}[id=TOAS04-P2-PHY-S3-Q04, source={TOAS - Test 04 (15 Mar 2022)}, paper=2, qno=4, type=numerical, subtopic={Torque}, answer={1.24}, answer_src=key, difficulty=5, diff_parts={K2 S2 I0 C1 F1}, diff_note={concepts: rotational energy, work by the field | steps: work of the electric field and gravity; moment of inertia; angular speed}, group=TOAS04-P2-PHY-S3-G3-4]
If electric field $E_0$ is switched on, the maximum angular speed acquired by the rod during its motion is $\sqrt{\dfrac{xg}{L}}$. Find $x$.
\end{bankq}
\begin{bankq}[id=TOAS06-P1-PHY-S2-Q10, source={TOAS - Test 06 (17 Apr 2022)}, paper=1, qno=10, type=multi, subtopic={Torque}, answer={A,C,D}, answer_src=key, difficulty=5, diff_parts={K2 S2 I1 C0 F1}, diff_note={concepts: rod equilibrium, spring energy | steps: energy as a function of theta; equilibrium condition; stability from the second derivative | I: energy method for the constrained rod}]
Two identical discs, each of negligible mass and radius $r(\lll\ell)$ are connected by a uniform rod of mass $m$ and length $\ell$ as shown in the figure. The disc A is connected with an ideal spring of spring constant $k$ at its centre. The spring is unstretched when $\theta=0^\circ$. There is no slipping at any contact at any point of time during the motion. Choose the correct statement(s) [Spring always obeys Hook's Law and $mg=k\,\ell$]. During the motion, each disc remains in contact with the contact surface :-
\bfig{132.5}{TOAS06_p1_phy_s2_q10_1.png}
\begin{qoptions}
\opt{A}{The value of $\theta$, when the system is in stable equilibrium is $60^\circ$}
\opt{B}{The value of $\theta$, when the system is in stable equilibrium is $30^\circ$}
\opt{C}{Spring potential energy of the system in the state of stable equilibrium is $\dfrac{k\ell^{2}}{8}$}
\opt{D}{The angular speed of the rod when it passes through the position of the stable equilibrium is $\sqrt{\dfrac{3g}{4\ell}}$}
\end{qoptions}
\end{bankq}
\begin{bankq}[id=TOAS09-P1-PHY-S1-Q01, source={TOAS - Test 09 (30 May 2022)}, paper=1, qno=1, type=single, subtopic={Moment of inertia}, answer={D}, answer_src=key, difficulty=1, diff_parts={K1 S1 I0 C0 F0}, diff_note={concepts: moment of inertia of a lamina | steps: I = m a\^{}2/12 for an equilateral triangle about the centroid, perpendicular axis}]
Moment of inertia of a uniform equilateral triangular lamina of mass 20 kg and side length 75 cm about an axis perpendicular to plane and passing through centroid (in kg--m$^2$) is :-
\begin{qoptions}
\opt{A}{$\dfrac{360}{96}$}
\opt{B}{$\dfrac{360}{32}$}
\opt{C}{$\dfrac{360}{192}$}
\opt{D}{$\dfrac{180}{192}$}
\end{qoptions}
\end{bankq}
\begin{bankpassage}[id=TOAS09-P1-PHY-S3-G1-2, source={TOAS - Test 09 (30 May 2022)}, questions={1-2}]
A uniform circular disc of mass $M=4$ kg and radius $R=25$ cm is suspended in a vertical plane and hinged at its centre `O'. It is free to rotate about a horizontal axis passing through its centre and perpendicular to its plane. A particle of mass $m=0.5$ kg moving horizontally with a velocity $u=20$ m/s collides with disc at point `P' and gets stuck to it as shown. ($g=10$ m/s$^2$)
\bfig{98.4}{TOAS09_p1_phy_s3_para1_2_1.png}
\end{bankpassage}
\begin{bankq}[id=TOAS09-P1-PHY-S3-Q01, source={TOAS - Test 09 (30 May 2022)}, paper=1, qno=1, type=numerical, subtopic={Angular momentum}, answer={8.00}, answer_src=key, difficulty=2, diff_parts={K1 S1 I0 C0 F1}, diff_note={concepts: angular momentum conservation in a collision | steps: angular momentum about the hinge before and after; angular velocity}, group=TOAS09-P1-PHY-S3-G1-2]
Find the angular velocity of the disc just after collision (in rad/sec)
\end{bankq}
\begin{bankq}[id=TOAS09-P1-PHY-S3-Q02, source={TOAS - Test 09 (30 May 2022)}, paper=1, qno=2, type=numerical, subtopic={Angular momentum}, answer={38.00}, answer_src=key, difficulty=4, diff_parts={K2 S2 I0 C0 F1}, diff_note={concepts: circular motion, hinge reaction | steps: energy conservation to the highest point; centripetal equation; force on the hinge}, group=TOAS09-P1-PHY-S3-G1-2]
The force on the disc due to hinge at its centre when the particle reaches the highest point (in newton).
\end{bankq}
\begin{bankq}[id=TOAS09-P2-PHY-S1-Q14, source={TOAS - Test 09 (30 May 2022)}, paper=2, qno=14, type=single, subtopic={Rotational kinematics}, answer={C}, answer_src=key, difficulty=1, diff_parts={K1 S1 I0 C0 F0}, diff_note={concepts: instantaneous axis of a rod | steps: direction of the velocity of A; angular velocity from the geometry}]
Instantaneous speed on end A of rod is `v'. Angular velocity of rod of length `L' will be.
\bfig{103.5}{TOAS09_p2_phy_s1_q14_1.png}
\begin{qoptions}
\opt{A}{$\dfrac{v}{2L}$}
\opt{B}{$\dfrac{v\sqrt3}{2L}$}
\opt{C}{$\dfrac{v}{L}$}
\opt{D}{$\dfrac{4v}{2L}$}
\end{qoptions}
\end{bankq}
\begin{bankq}[id=TOAS09-P2-PHY-S1-Q15, source={TOAS - Test 09 (30 May 2022)}, paper=2, qno=15, type=single, subtopic={Torque}, answer={A}, answer_src=key, difficulty=1, diff_parts={K1 S1 I0 C0 F0}, diff_note={concepts: toppling of a prism | steps: torque about the pivot edge; force for toppling}]
Equilateral prism of mass $m$ is rest on rough horizontal floor. If coefficient of friction is sufficiently high such that prism does not slide before topple. Minimum force required to topple prism is (coefficient is friction = $\mu$)
\bfig{77.7}{TOAS09_p2_phy_s1_q15_A.png}
\begin{qoptions}
\opt{A}{$\dfrac{mg}{\sqrt3}$}
\opt{B}{$\dfrac{\mu mg}{\sqrt3}$}
\opt{C}{$\dfrac{mg}{4}$}
\opt{D}{$\dfrac{\mu mg}{4}$}
\end{qoptions}
\end{bankq}
\begin{bankq}[id=TOAS09-P2-PHY-S1-Q16, source={TOAS - Test 09 (30 May 2022)}, paper=2, qno=16, type=single, subtopic={Rolling}, answer={D}, answer_src=key, difficulty=1, diff_parts={K1 S1 I0 C0 F0}, diff_note={concepts: rolling disc under a tangential force | steps: torque about the centre and rolling constraint; acceleration of the centre}]
A force `F' act tangentially at highest point of a disc of mass `m' \& radius `R'. If disc is rolling without sliding. Acceleration of centre of disc will be
\bfig{109.5}{TOAS09_p2_phy_s1_q16_1.png}
\begin{qoptions}
\opt{A}{$\dfrac{2F}{3m}$}
\opt{B}{0}
\opt{C}{$\dfrac{10F}{7m}$}
\opt{D}{$\dfrac{4F}{3m}$}
\end{qoptions}
\end{bankq}
\begin{bankq}[id=TOAS09-P2-PHY-S2-Q07, source={TOAS - Test 09 (30 May 2022)}, paper=2, qno=7, type=numerical, subtopic={Moment of inertia}, answer={2}, answer_src=solved, working={Base rod has zero moment; each slanted rod $\frac13m\ell^2\sin^2 60^\circ=\frac{m\ell^2}4$: total $\frac{m\ell^2}2$, $x=2$}, difficulty=3, diff_parts={K1 S2 I0 C0 F1}, diff_note={concepts: moment of inertia of rods about an axis | steps: zero for the base rod; slanted rods with distance; sum}]
Three identical rod each of mass $m$ \& length $\ell$ form a equilateral triangle. Moment of inertia about given axis is $\dfrac{m\ell^2}{x}$, value of x will be :-
\bfig{134.5}{TOAS09_p2_phy_s2_q07_1.png}
\end{bankq}
\begin{bankq}[id=TOAS09-P2-PHY-S2-Q08, source={TOAS - Test 09 (30 May 2022)}, paper=2, qno=8, type=numerical, subtopic={Torque}, answer={1500}, answer_src=solved, working={Equal tensions $T=(500+W)/2$; torque about the left end: $1.3T=250+0.7W$, $W=1500$ N}, difficulty=2, diff_parts={K1 S1 I0 C0 F1}, diff_note={concepts: equilibrium of a plank | steps: force balance with equal tensions; torque balance gives W}]
In given diagram a uniform plank weight is 500 N. Find out weight W if $T_1$ \& $T_2$ are equal.
\bfig{123.9}{TOAS09_p2_phy_s2_q08_1.png}
\end{bankq}
\begin{bankq}[id=TOAS10-P1-PHY-S1-Q02, source={TOAS - Test 10 (12 Jun 2022)}, paper=1, qno=2, type=multi, subtopic={Rotational kinematics}, answer={A,B,C,D}, answer_src=key, difficulty=5, diff_parts={K2 S2 I1 C0 F1}, diff_note={concepts: rigid body kinematics, energy conservation | steps: no horizontal force so the centre of mass falls on a vertical line; instantaneous centre and energy give omega at the horizontal position; acceleration of the centre of mass from the normal reaction; locus of an interior point | I: instantaneous centre of rotation}]
A thin rod of length $\ell$ is kept vertically (along y-axis) with end A being at rest on a smooth horizontal plane. A slight disturbance makes the end A slip on the plane (along x-axis) and the rod falls. Then,
\begin{qoptions}
\opt{A}{the path described by a point P at a distance of $\dfrac{\ell}{4}$ from A is a quadrant of a circle.}
\opt{B}{the centre of mass of the rod moves along a straight line, parallel to y-axis.}
\opt{C}{the acceleration of the centre of mass as the rod becomes horizontal is $\dfrac{3g}{4}$.}
\opt{D}{the greatest angular velocity of the rod about the instantaneous centre of rotation is $\sqrt{\dfrac{3g}{\ell}}$.}
\end{qoptions}
\end{bankq}
\begin{bankq}[id=TOAS10-P1-PHY-S1-Q03, source={TOAS - Test 10 (12 Jun 2022)}, paper=1, qno=3, type=multi, subtopic={Moment of inertia}, answer={A,D}, answer_src=key, difficulty=3, diff_parts={K1 S2 I0 C0 F1}, diff_note={concepts: moment of inertia with parallel axis theorem | steps: lamina about the centroid; particles about the centroid; shift to A by parallel axes; shift to C}]
A, B and C are the vertices of a uniform equilateral triangular thin lamina of mass M. Three particle of masses $\dfrac M3$, each are located at the midpoints P, Q and R of the edges of the lamina. For this system, (each edge of lamina = a)
\bfig{84.6}{TOAS10_p1_phy_s1_q03_1.png}
\begin{qoptions}
\opt{A}{the moment of inertia of the system about an axis through the centroid O perpendicular to the plane is $\dfrac{Ma^2}{6}$.}
\opt{B}{the moment of inertia of the system about an axis through O perpendicular to the plane is $\dfrac{5Ma^2}{12}$}
\opt{C}{the moment of inertia of the system about an axis through A perpendicular to the plane is $\dfrac{13Ma^2}{12}$}
\opt{D}{the moment of inertia of the system about an axis through C perpendicular to the plane is $\dfrac{5Ma^2}{6}$}
\end{qoptions}
\end{bankq}
\begin{bankq}[id=TOAS10-P1-PHY-S1-Q07, source={TOAS - Test 10 (12 Jun 2022)}, paper=1, qno=7, type=multi, subtopic={Torque}, answer={A,B,C,D}, answer_src=key, difficulty=2, diff_parts={K1 S1 I0 C0 F1}, diff_note={concepts: rotation of a disc with a hanging mass | steps: acceleration and tension from the two equations; angular displacement, torque work and kinetic energy in 4 s}]
A uniform disc of mass 2kg and radius 1m is mounted on a horizontal axle supported on fixed frictionless bearings. A light chord is wrapped around the rim of the disc and mass of 1kg is tied to the free end. If it is released from rest (take $g=10$ m/s$^2$)
\begin{qoptions}
\opt{A}{the tension in the chord is 5N.}
\opt{B}{in first four seconds the angular displacement of the disc is 40 rad.}
\opt{C}{the work done by the torque on the disc in first four seconds is 200 J.}
\opt{D}{the increase in the kinetic energy of the mass in the first four seconds is 200 J.}
\end{qoptions}
\end{bankq}
\begin{bankq}[id=TOAS10-P1-PHY-S1-Q08, source={TOAS - Test 10 (12 Jun 2022)}, paper=1, qno=8, type=multi, subtopic={Angular momentum}, answer={C,D}, answer_src=key, difficulty=5, diff_parts={K2 S2 I1 C0 F1}, diff_note={concepts: angular momentum about the edge, restitution | steps: angular momentum about the edge; restitution at the contact; velocity of the centre of mass; check the statements | I: impulse acts only at the contact point}]
A uniform rod of length $\ell$ is falling down with a velocity $v_0$ when one of it's end hits a fixed edge as shown. (Consider impulsive force to be acting perpendicular to the length of rod)
\bfig{93.1}{TOAS10_p1_phy_s1_q08_1.png}
\begin{qoptions}
\opt{A}{If the collision is elastic, the centre of mass may have upwards velocity just after the collision.}
\opt{B}{If the collision is inelastic, the centre of mass may come to rest just after the collision.}
\opt{C}{If the collision is perfectly inelastic, the centre of mass has a downward velocity just after the collision.}
\opt{D}{If the collision is elastic, the centre of mass has a downward velocity just after the collision.}
\end{qoptions}
\end{bankq}
\begin{bankpassage}[id=TOAS10-P1-PHY-S2-G9-10, source={TOAS - Test 10 (12 Jun 2022)}, questions={9-10}]
A small particle of mass $m$ is fixed to the perimeter of a ring of same mass and radius `r' and ring is placed on a horizontal plane. Friction is negligible on horizontal plane. Initially particle is at top most point, then the system is released from rest. Answer next two questions when the particle is at the same height as the centre of ring after being released from topmost point. Assume that the ring stays in vertical plane during its motion under consideration.
\bfig{93.3}{TOAS10_p1_phy_s2_para9_10_1.png}
\end{bankpassage}
\begin{bankq}[id=TOAS10-P1-PHY-S2-Q09, source={TOAS - Test 10 (12 Jun 2022)}, paper=1, qno=9, type=multi, subtopic={Rotational kinematics}, answer={A,B,D}, answer_src=key, difficulty=5, diff_parts={K2 S2 I1 C0 F1}, diff_note={concepts: motion of ring, particle on a smooth plane | steps: no horizontal force so the horizontal momentum is zero; energy conservation for the system; contact point and centre velocities | I: horizontal momentum zero}, group=TOAS10-P1-PHY-S2-G9-10]
Mark the correct option(s) :-
\begin{qoptions}
\opt{A}{Speed of centre of mass of system is $\sqrt{\dfrac{gr}{4}}$.}
\opt{B}{Angular velocity of system is $\omega=\sqrt{\dfrac gr}$}
\opt{C}{Contact point of ring with ground is instantaneously at rest.}
\opt{D}{Centre of ring is instantaneously at rest.}
\end{qoptions}
\end{bankq}
\begin{bankq}[id=TOAS10-P1-PHY-S2-Q10, source={TOAS - Test 10 (12 Jun 2022)}, paper=1, qno=10, type=multi, subtopic={Rotational kinematics}, answer={A,C}, answer_src=key, difficulty=6, diff_parts={K2 S2 I1 C1 F1}, diff_note={concepts: rigid body dynamics, centre of mass motion | steps: instantaneous acceleration of the particle and the centre; angular acceleration; normal force from the vertical equation; direction of acceleration | I: use the centre of mass motion}, group=TOAS10-P1-PHY-S2-G9-10]
Mark the correct option(s) :
\begin{qoptions}
\opt{A}{Instantaneous acceleration of the centre of ring is $\dfrac g2$.}
\opt{B}{Acceleration of centre of ring initially points towards left and then towards right.}
\opt{C}{The force exerted by the ring on the ground at the instant given is $\dfrac32mg$.}
\opt{D}{The force exerted by the ring on the ground at the instant given is $2mg$.}
\end{qoptions}
\end{bankq}
\begin{bankq}[id=TOAS10-P1-PHY-S3-Q08, source={TOAS - Test 10 (12 Jun 2022)}, paper=1, qno=8, type=numerical, subtopic={Rolling}, answer={3.00}, answer_src=key, difficulty=4, diff_parts={K1 S2 I1 C0 F1}, diff_note={concepts: rolling with slipping | steps: velocity of a rim point; condition that its horizontal velocity is zero; height above the ground | I: slipping disc has no rolling constraint}]
A disc of radius `5cm' rolls on a horizontal surface with linear velocity $v=1\hat i$ m/s and angular velocity 50 rad/sec. Height of particle from ground on rim of disc which has velocity in vertical direction is (in cm) :
\bfig{131.6}{TOAS10_p1_phy_s3_q08_1.png}
\end{bankq}
\begin{bankq}[id=TOAS12-P2-PHY-S1-Q05, source={TOAS - Test 12 (10 Jul 2022)}, paper=2, qno=5, type=single, subtopic={Torque}, answer={C}, answer_src=key, difficulty=4, diff_parts={K2 S2 I1 C0 F0}, diff_note={concepts: rotation of a rod, electric force | steps: torque of electric and gravitational forces; condition that the rod just becomes horizontal; maximum angular speed at equilibrium by energy | I: maximum speed at the torque-free position}]
A uniformly charged non-conducting rod is suspended vertically at its end. The rod can swing freely in the vertical plane without any friction. The linear charge density on the rod is $\lambda$ C/m and it has a uniformly distributed mass $M$. Length of the rod is $L$. A uniform horizontal field $E$ is switched on so that the rod just manages to make itself horizontal. What is the maximum angular speed of the rod during its motion ?
\begin{qoptions}
\opt{A}{$\sqrt{\dfrac{3(\sqrt3-1)g}{L}}$}
\opt{B}{$\sqrt{\dfrac{3(\sqrt2-1)g}{2L}}$}
\opt{C}{$\sqrt{\dfrac{3(\sqrt2-1)g}{L}}$}
\opt{D}{$\sqrt{\dfrac{2(\sqrt3-1)g}{L}}$}
\end{qoptions}
\end{bankq}
\begin{bankq}[id=TOAS15-P2-PHY-S2-Q01, source={TOAS - Test 15 (22 Aug 2022)}, paper=2, qno=1, type=numerical, subtopic={Moment of inertia}, answer={0.37}, answer_src=key, difficulty=4, diff_parts={K2 S2 I0 C0 F1}, diff_note={concepts: torsional oscillation, moment of inertia | steps: frequency scales as 1/sqrt(I); added moment 2 m (L/2)\^{}2; ratio from the 20 percent drop}]
A rod of mass `$M$' and length `$2L$' is suspended at its middle by a wire. It exhibits torsional oscillations; If two masses each of `$m$' are attached at distance `$L/2$' from its centre on both sides, it reduces the oscillation frequency by 20\%. The value of ratio $m/M$ is close to:
\end{bankq}
\begin{bankq}[id=TOAS19-P1-PHY-S1-Q05, source={TOAS - Test 19 (05 Feb 2023)}, paper=1, qno=5, type=multi, subtopic={Torque}, answer={A,B,C,D}, answer_src=key, difficulty=3, diff_parts={K1 S2 I0 C0 F1}, diff_note={concepts: equilibrium of a three-legged table | steps: total weight; extra loads sit on the legs; share of the board weight; effect of removal}]
Figure shows three leg table, with top as a uniform equilateral triangular board of mass 10 kg. Two bodies of masses 2 kg and 3 kg are placed at corners A and C. The table is placed on a horizontal surface.
\bfig{74.0}{TOAS19_p1_phy_s1_q05_1.png}
\begin{qoptions}
\opt{A}{Normal force on leg at A $=\dfrac{160}3$ N}
\opt{B}{Normal force on leg at B $=\dfrac{100}3$ N}
\opt{C}{Normal force on leg at C $=\dfrac{190}3$ N}
\opt{D}{If both the blocks are removed, normal force on leg at B will remain same.}
\end{qoptions}
\end{bankq}
\begin{bankq}[id=TOAS19-P1-PHY-S1-Q08, source={TOAS - Test 19 (05 Feb 2023)}, paper=1, qno=8, type=multi, subtopic={Moment of inertia}, answer={A,C}, answer_src=key, difficulty=8, diff_parts={K2 S3 I2 C1 F1}, diff_note={concepts: rigid rod in a circular groove, rotation about the centre | steps: rod rotates rigidly about the centre; moment of inertia; torque and acceleration of the centre of mass; energy for the maximum speed; normal force when horizontal}]
Figure shows a smooth circular groove in vertical plane in which a uniform rod of length $\ell=1.2\,R$ is released from the position shown.
\bfig{89.9}{TOAS19_p1_phy_s1_q08_1.png}
\begin{qoptions}
\opt{A}{Initial acceleration of centre of mass of rod is $\dfrac{16}{19}g$.}
\opt{B}{Maximum speed of centre of mass of rod will be $\sqrt{1.6gR}$}
\opt{C}{Maximum angular speed of rod will be $\sqrt{\dfrac{40}{19}\dfrac gR}$}
\opt{D}{Normal force at one end of rod when it becomes horizontal is $mg$.}
\end{qoptions}
\end{bankq}
\begin{bankq}[id=TOAS19-P1-PHY-S2-Q09, source={TOAS - Test 19 (05 Feb 2023)}, paper=1, qno=9, type=single, subtopic={Torque}, answer={A}, answer_src=key, difficulty=5, diff_parts={K2 S2 I1 C1 F0}, diff_note={concepts: pure rolling wheelie, torque balance | steps: friction from M a; engine torque from the wheel equation; torque balance about the centre of mass; h and b}]
Figure shows a rider performing a stunt with rear wheel in pure rolling.

$M=$ total mass of system $=200$ kg,

$m=$ mass of rear wheel (ring) $=10$ kg,

$R=$ radius of wheel $=50$ cm,

$h=$ height of centre of mass of system,

$b=$ distance of vertical line through centre of mass from contact point of rear wheel.
\bfig{150.8}{TOAS19_p1_phy_s2_q09_1.png}
\begin{center}\footnotesize\setlength{\tabcolsep}{2pt}
\begin{tabular}{|L{0.8cm}L{3.6cm}|L{0.6cm}L{1.0cm}|}\hline
\multicolumn{2}{|c|}{List-I}&\multicolumn{2}{c|}{List-II}\\\hline
(I)&Frictional force in KN&(P)&1.2\\\hline
(II)&Torque provided by engine in KN-m&(Q)&0.630\\\hline
(III)&Value of $h$ in meter&(R)&50\\\hline
(IV)&Value of $b$ in centimeter&(S)&1\\\hline
&&(T)&60\\\hline
&&(U)&0.525\\\hline
\end{tabular}\end{center}
It the motorcycle in the given situation accelerates with 5 m/s$^2$, the correct matching is as :-
\begin{qoptions}
\opt{A}{I $\to$ S;II $\to$ U;III $\to$ S;IV $\to$ R}
\opt{B}{I $\to$ S;II $\to$ U;III $\to$ R;IV $\to$ S}
\opt{C}{I $\to$ S;II $\to$ Q;III $\to$ S;IV $\to$ S}
\opt{D}{I $\to$ P;II $\to$ Q;III $\to$ R;IV $\to$ S}
\end{qoptions}
\end{bankq}
\begin{bankq}[id=TOAS19-P1-PHY-S2-Q10, source={TOAS - Test 19 (05 Feb 2023)}, paper=1, qno=10, type=single, subtopic={Torque}, answer={A}, answer_src=key, difficulty=5, diff_parts={K2 S2 I1 C1 F0}, diff_note={concepts: pure rolling wheelie, torque balance | steps: friction from M a; engine torque including the ring inertia; h and b for the new acceleration}]
Figure shows a rider performing a stunt with rear wheel in pure rolling.

$M=$ total man of system $=200$ kg,

$m=$ mass of rear wheel (ring) $=10$ kg,

$R=$ radius of wheel $=50$ cm,

$h=$ height of centre of mass of system,

$b=$ distance of vertical line through centre of mass from contact point of rear wheel.
\bfig{150.8}{TOAS19_p1_phy_s2_q10_1.png}
\begin{center}\footnotesize\setlength{\tabcolsep}{2pt}
\begin{tabular}{|L{0.8cm}L{3.6cm}|L{0.6cm}L{1.0cm}|}\hline
\multicolumn{2}{|c|}{List-I}&\multicolumn{2}{c|}{List-II}\\\hline
(I)&Frictional force in KN&(P)&1.2\\\hline
(II)&Torque provided by engine in KN-m&(Q)&0.630\\\hline
(III)&Value of $h$ in meter&(R)&50\\\hline
(IV)&Value of $b$ in centimeter&(S)&1\\\hline
&&(T)&60\\\hline
&&(U)&0.525\\\hline
\end{tabular}\end{center}
It the motorcycle in the given situation accelerates with 6 m/s$^2$, the correct matching is as :-
\begin{qoptions}
\opt{A}{I $\to$ P;II $\to$ Q;III $\to$ S;IV $\to$ T}
\opt{B}{I $\to$ P;II $\to$ Q;III $\to$ T;IV $\to$ S}
\opt{C}{I $\to$ P;II $\to$ S;III $\to$ R;IV $\to$ Q}
\opt{D}{I $\to$ R;II $\to$ S;III $\to$ T;IV $\to$ U}
\end{qoptions}
\end{bankq}
\end{document}
